\documentclass[12pt]{article}
\usepackage[utf8]{inputenc}
\usepackage[T1]{fontenc}
\usepackage[margin=1in]{geometry}
\usepackage{setspace}
\usepackage{xcolor}
\usepackage[colorlinks=true, urlcolor=blue!60!black, linkcolor=blue!60!black]{hyperref}
\setlength{\parindent}{0em}
\setlength{\parskip}{0.6em}
\renewcommand{\baselinestretch}{1.0}
\begin{document}
\begin{center}
{\LARGE \textit{South Dakota v. Dole}, 483 U.S. 203 (1987)\par}
\vspace{0.5em}
{\large Case Brief --- Two-Page Edition\par}
\end{center}

\section{Memory Jogger}
Upholding a law withholding 5\% of highway funds from States with drinking ages under 21, the Court set the conditional-spending test: general welfare, clear notice, relatedness, no independent constitutional bar, and no coercion (\href{https://supreme.justia.com/cases/federal/us/483/203/}{483 U.S. 203}).

\section{Detailed Case Facts}
To stop young adults from crossing state lines to drink and then drive, Congress in 1984 directed the Secretary of Transportation to withhold 5\% of specified highway funds from any State allowing people under 21 to buy or publicly possess alcohol. South Dakota let 19-year-olds buy low-alcohol beer.

\section{Procedural History}
South Dakota sued Secretary Elizabeth Dole under the Spending Clause and the Twenty-first Amendment. The District Court and the Eighth Circuit rejected its claims.

\section{Judicial Votes}
Affirmed, 7--2. Rehnquist wrote for White, Marshall, Blackmun, Powell, Stevens, and Scalia; Brennan and O'Connor dissented separately.

\section{Holding}
Even assuming Congress could not set a drinking age directly, it may induce States to do so through conditional grants that serve the general welfare, are unambiguous, relate to the federal interest, require nothing unconstitutional, and stop short of coercion.

\section{Analysis of Opinions}
Following \textit{Butler}, Rehnquist held that the spending power reaches beyond Article I's enumerated fields. The condition was clear and ``directly related'' to safe interstate travel. The independent-bar limit forbids only inducing unconstitutional conduct, and raising the drinking age was lawful. Losing 5\% of highway funds was ``relatively mild encouragement.'' Brennan found the drinking age reserved to the States. O'Connor found the condition ``far too over- and under-inclusive'' and argued that Congress may dictate how its money is spent but not regulate who may drink.

\section{Examples: Future Applications}
\textbf{Same-Side:} withholding 5\% of highway funds from States without a 0.08 blood-alcohol limit is valid. \textbf{Opposite-Side:} threatening all existing K--12 funds to force a campus gun ban is ``a gun to the head'' under \textit{NFIB}. \textbf{Fence-Sitter:} withholding 25\% of education funds unless States ban student phones.

\section{Critique}
Scholars see a ``Trojan Horse'' letting Congress buy compliance it could not mandate, ironically from the author of \textit{Lopez}. Neither \textit{Dole} nor \textit{NFIB} offers a bright line for relatedness or coercion.

\section{Key Quotations}
\begin{quote}
``These cases establish that the `independent constitutional bar' limitation on the spending power is not, as petitioner suggests, a prohibition on the indirect achievement of objectives which Congress is not empowered to achieve directly. Instead, we think that the language in our earlier opinions stands for the unexceptionable proposition that the power may not be used to induce the States to engage in activities that would themselves be unconstitutional.''
\end{quote}
\noindent --- Chief Justice Rehnquist, Opinion of the Court, 483 U.S. at 210.

\end{document}
